% Cote 119: offprints, reports, courses, CV and letters (typed and printed part): archive cote 119, site page 254 (PDF page 255).
% typescript. Draft compared: none.
% AI-transcribed by Claude Opus 5.5 (claude-opus-5-5, adaptive thinking, effort high), Message Batches API, batch r119_full (msgbatch_01DwZNm5Y7YsKBWyAN4sFs8w), 2026-11-13; not yet checked by eye.
\scanpage{119/254}
\margin{\unsure{4,7,1,}} % NOTE: at top right a dark mark (staple or tear) hides part of a handwritten number, of which ``4,7,1,'' and a final ``7'' are visible
\margin{7}

Le maître-enseignant et le maître-chercheur dans
l'Université d'aujourd'hui et de demain

par A. Grothendieck

A. \uline{Aujourd'hui}

1. Dans le système universitaire tel qu'il a existé jusqu'à nos jours, depuis
\margin{|} semble-t-il ses origines lointaines, le rôle du maître, d'une part comme l'en-%
seignant chargé de transmettre des connaissances acquises, et d'autre part
comme créateur intellectuel contribuant à l'approfondissement, à l'élargis-%
sement et au renouvellement de ces mêmes connaissances, était confondu. En
d'autres termes, une même personne était censée remplir ces deux rôles, qui
sont en effet étroitement apparentés, mais néanmoins essentiellement distincts.
N'était censé être qualifié pour enseigner la Science que celui qui pouvait
prétendre faire partie de ceux qui contribuaient à l'édifier.

Ce système était effectivement compatible avec les besoins de la
société et de la science pendant fort longtemps. Les études universitaires
étaient réservées à un petit nombre, et les étudiants dans chaque discipline
poursuivaient généralement leurs études à un niveau théorique élevé, qui les
menait jusqu'aux frontières mouvantes entre le connu et le domaine de la
recherche. Il était naturel que l'enseignement universitaire soit assuré par
des maîtres dont chacun dominait à peu près complètement la science qu'il
enseignait, et dont chacun était lui-même un des artisans de cette science.
Vu le nombre limité d'étudiants, des effectifs modestes suffisaient pour les
enseigner, et le nombre des savants de valeur susceptibles d'assurer cet
\margin{|} enseignement était largement suffisant. Il n'était pas rare, même, que des
savants éminents soient obligés d'attendre de longues années la retraite d'un
de leurs collègues pour pouvoir occuper une position de titulaire de chaire
à la mesure de leur valeur.

2. Avec l'essor de la technologie dans ces dernières cinquante années, et
l'afflux de plus en plus considérable d'étudiants dans les Facultés, et notam-%
ment dans les Facultés de Sciences, cette situation a radicalement changé, non

\mspage{254}
